\chapter{Compartment syndrome}
\index{Compartment syndrome}

\section*{Definition and Overview}
Compartment syndrome is a potentially limb-threatening and life-threatening medical emergency characterised by increased pressure within a closed anatomical osteofascial compartment. This elevation in tissue pressure compromises local microvascular circulation, leading to tissue ischaemia, cellular hypoxia, and progressive neuromuscular necrosis. Compartment syndrome most commonly occurs in the lower leg and forearm but can develop in any enclosed space, including the thigh, foot, hand, and abdomen. It is broadly categorised into acute compartment syndrome (ACS), which requires immediate surgical decompression via fasciotomy, and chronic exertional compartment syndrome (CECS), which is a non-emergent condition associated with exercise-induced pain.

\section*{Epidemiology}
Acute compartment syndrome can occur at any age but is most prevalent in young males under 35 years of age, which correlates with higher rates of high-energy musculoskeletal trauma. The most common predisposing cause is a fracture of a long bone, accounting for approximately 75\% of ACS cases, with tibial shaft fractures representing the single highest risk (up to 10\% of such fractures develop ACS), followed by fractures of the radius and ulna. Non-traumatic causes, such as crush injuries, drug overdoses with prolonged limb compression, and reperfusion injury, are less common but significant. Chronic exertional compartment syndrome typically affects endurance athletes, runners, and military recruits, with a high proportion of bilateral involvement.

\section*{Aetiology and Pathophysiology}
The pathophysiology of compartment syndrome is driven by a mismatch between the volume of the osteofascial compartment and its contents.
\begin{itemize}
    \item \textbf{Decreased compartment size:} Caused by external compression, such as tight casts, splints, dressings, or circumferential burn eschars.
    \item \textbf{Increased compartment contents:} Caused by internal swelling or fluid accumulation, including haemorrhage (due to fractures or vascular injury), oedema (from ischaemia-reperfusion injury, trauma, or venomous bites), and capillary leak.
    \item \textbf{Pressure dynamics:} Muscle groups are enclosed by rigid fascial envelopes. As tissue volume expands, pressure rises. When intracompartmental pressure exceeds capillary perfusion pressure (typically within 10 to 30 mmHg of the diastolic blood pressure), capillaries collapse.
    \item \textbf{Ischaemic cascade:} Capillary collapse halts microvascular perfusion, leading to cellular hypoxia. Muscle tissue can tolerate up to 4 hours of ischaemia before necrosis begins, whereas peripheral nerves suffer irreversible damage after approximately 6 to 8 hours.
    \item \textbf{Reperfusion injury:} Restoring blood flow to ischaemic muscle can paradoxically exacerbate swelling due to oxygen free radical release, capillary leak, and massive cell lysis.
\end{itemize}

\section*{Clinical Features}
The clinical diagnosis of acute compartment syndrome relies on the classic "5 Ps", although many of these are late signs of irreversible tissue damage.
\begin{itemize}
    \item \textbf{Pain out of proportion:} The earliest and most reliable sign; pain is severe, progressive, and poorly responsive to opioid analgesics.
    \item \textbf{Pain with passive stretch:} Exquisite pain elicited when the muscles of the affected compartment are passively stretched (e.g.\ passive extension of the toes in anterior tibial compartment syndrome).
    \item \textbf{Paresthesia:} Sensory deficit or altered sensation (tingling, numbness) in the distribution of the nerves passing through the compartment, indicating early nerve ischaemia.
    \item \textbf{Pressure and fullness:} The affected limb feels palpably tense, hard, swollen, and wood-like on physical examination.
    \item \textbf{Pallor and pulselessness:} Late and ominous signs indicating major arterial occlusion and impending limb loss; the presence of distal pulses does \textbf{not} rule out compartment syndrome.
    \item \textbf{Paralysis:} Late sign representing muscle necrosis and nerve injury, resulting in inability to actively contract the affected muscles.
\end{itemize}

\section*{Differential Diagnosis}
It is critical to distinguish acute compartment syndrome from other causes of acute limb pain and swelling to ensure prompt intervention:
\begin{itemize}
    \item \textbf{Deep vein thrombosis (DVT):} Presents with unilateral limb swelling, warmth, and calf pain, but lacks pain with passive stretch and paresthesia, and pressure is not localized to a single fascial compartment.
    \item \textbf{Cellulitis or necrotising fasciitis:} Infectious conditions presenting with erythema, warmth, and systemic symptoms (fever, leucocytosis); pain in necrotising fasciitis is also severe but is accompanied by soft-tissue crepitus and skin necrosis.
    \item \textbf{Arterial occlusion:} Acute peripheral arterial embolism presents with sudden pain, pulselessness, and coldness, but lacks the localized compartmental swelling and hardness.
    \item \textbf{Osteomyelitis:} Bone infection causing localized pain and swelling, but typically presents with systemic signs of infection and lack of acute compartmental pressure rises.
    \item \textbf{Shin splints (medial tibial stress syndrome):} A chronic cause of exercise-induced leg pain that resolves quickly with rest, whereas chronic exertional compartment syndrome persists longer and presents with transient sensory changes.
\end{itemize}

\section*{When to Seek Medical Care}
Acute compartment syndrome is a surgical emergency. Any patient who has recently sustained a limb injury, fracture, or has a tight cast, and develops worsening, unmanageable pain must seek emergency medical care immediately.

\textbf{Contact your local emergency services for:}
\begin{itemize}
    \item Extreme, progressive pain in a limb that is out of proportion to the injury or is not relieved by strong painkillers.
    \item Numbness, tingling, or loss of sensation in the foot, toes, hand, or fingers after an injury.
    \item A cold, pale, or bluish hand or foot, particularly if pulses are weak or absent.
    \item Inability to move the fingers or toes.
    \item A feeling of extreme tightness, hardness, or swelling in the arm or leg.
\end{itemize}

\section*{Diagnosis and Investigations}
The diagnosis of acute compartment syndrome is primarily clinical. Objective measurements of intracompartmental pressure are used when the clinical picture is equivocal or the patient is obtunded.
\begin{itemize}
    \item \textbf{Intracompartmental pressure (ICP) monitoring:} Measured using a needle or catheter attached to a digital monitor inserted directly into the affected compartment. An ICP greater than 30 mmHg, or a delta pressure ($\Delta P$, calculated as diastolic blood pressure minus ICP) of less than 30 mmHg, is diagnostic.
    \item \textbf{Serum creatine kinase (CK):} Measured to assess the degree of muscle damage; levels are elevated in muscle necrosis and rhabdomyolysis.
    \item \textbf{Urine myoglobin:} Detected in cases of severe muscle necrosis to monitor for the development of pigment-induced acute kidney injury.
    \item \textbf{Radiography (X-rays):} Performed immediately to identify underlying fractures or dislocations that require urgent reduction.
    \item \textbf{Duplex ultrasonography:} May be used to rule out venous thrombosis, though it should not delay definitive surgical treatment if compartment syndrome is suspected.
\end{itemize}

\section*{Management}
The management of acute compartment syndrome requires urgent decompression of the affected compartments.
\begin{itemize}
    \item \textbf{Immediate conservative measures:} Elevate the limb to the level of the heart (not above, as hyper-elevation decreases arterial pressure and perfusion), remove all tight casts, bandages, and splints, and provide oxygen and intravenous fluids to optimise blood flow.
    \item \textbf{Emergency surgical fasciotomy:} The definitive treatment. Long surgical incisions are made through the skin and fascia to release pressure in all affected compartments of the limb.
    \item \textbf{Wound management:} Fasciotomy wounds are left open and covered with sterile dressings or negative pressure wound therapy (NPWT) to allow swelling to subside.
    \item \textbf{Delayed primary closure:} Wounds are closed or skin grafted 3 to 5 days later once swelling has resolved and the tissue is viable.
    \item \textbf{Treatment of complications:} Intravenous hydration and urine alkalisation to prevent acute kidney injury from rhabdomyolysis.
    \item \textbf{Physical rehabilitation:} Aggressive physiotherapy to restore joint range of motion, muscle strength, and prevent contractures (such as Volkmann's ischaemic contracture).
\end{itemize}

\section*{Complications}
\begin{itemize}
    \item \textbf{Muscle necrosis and fibrosis:} Irreversible death of muscle tissue, leading to loss of function and replacing muscle with scar tissue.
    \item \textbf{Volkmann's ischaemic contracture:} Permanent flexion contracture of the wrist and hand (or ankle and foot) resulting from ischaemic muscle necrosis.
    \item \textbf{Permanent nerve damage:} Chronic sensory loss and motor deficits (e.g.\ foot drop from common peroneal nerve injury).
    \item \textbf{Infection and osteomyelitis:} High risk of bacterial contamination and wound infection in open fasciotomy wounds, which can spread to underlying bones.
    \item \textbf{Rhabdomyolysis and renal failure:} Release of myoglobin from necrotic muscle cells into the bloodstream, blocking renal tubules and causing acute kidney injury.
    \item \textbf{Amputation:} Indicated if decompression is delayed and the limb is non-viable or severely infected.
\end{itemize}

\section*{Prognosis and Outcomes}
The prognosis of acute compartment syndrome is heavily dependent on the time elapsed between symptom onset and surgical decompression. If fasciotomy is performed within 6 hours of onset, excellent recovery of muscle and nerve function is achieved in over 90\% of cases. If surgery is delayed beyond 12 hours, permanent neuromuscular disability is common, and rates of infection and amputation rise significantly. Chronic exertional compartment syndrome carries a highly favourable prognosis, with most patients achieving symptom resolution through activity modification or elective, minimally invasive surgical release.

\section*{Prevention and Self-Care}
\begin{itemize}
    \item \textbf{Monitor tight casts:} Check limbs in splints or casts frequently for swelling, discoloration, coldness, or numbness.
    \item \textbf{Educate post-injury patients:} Inform patients discharged with fractures of the signs of compartment syndrome and the importance of returning to the hospital immediately if pain worsens.
    \item \textbf{Prompt cast bivalving:} Splitting or removing casts immediately if a patient complains of worsening pain.
    \item \textbf{Gradual training progression:} For athletes, increasing exercise intensity slowly and incorporating rest days to prevent chronic exertional compartment syndrome.
    \item \textbf{Proper positioning:} Ensure injured limbs are supported at heart level to maintain perfusion while waiting for medical evaluation.
\end{itemize}

\section*{Sources and Further Reading}
The following reputable organisations provide further information on compartment syndrome:
\begin{itemize}
    \item Healthdirect Australia. \textit{Overview of compartment syndrome causes and treatment.} https://www.healthdirect.gov.au/
    \item OrthoInfo (American Academy of Orthopaedic Surgeons). \textit{Detailed guide to compartment syndrome.} https://orthoinfo.aaos.org/
    \item Melbourne Orthopaedic Group. \textit{Understanding chronic exertional compartment syndrome.} https://www.melbourneorthopaedicgroup.com.au/
    \item NHS. \textit{Compartment syndrome symptoms and emergency care.} https://www.nhs.uk/
    \item MSD Manual (Professional Edition). \textit{Pathophysiology and management of compartment syndrome.} https://www.msdmanuals.com/
\end{itemize}
